\chapter{Experiments and Results}
\label{chap:results}

\section{Introduction}

This chapter reports the fourteen training runs described in Chapter~\ref{chap:methodology} and
what they say about the research question. Section~\ref{sec:res-integrity} reports the checks
that decide whether the numbers can be trusted at all. Section~\ref{sec:res-repro} compares the
eight loss functions that the BAREC benchmark also reports, which is the evidence that the
pipeline is faithful. Section~\ref{sec:res-main} presents the full results on both splits and
answers the question put by the research brief: how CORAL and the QWK-oriented objective
compare with the distance-aware losses under one identical protocol.
Sections~\ref{sec:res-tradeoff} and~\ref{sec:res-distance} separate two things that the single
headline metric hides, namely exact-match accuracy and the distribution of error distances.
Sections~\ref{sec:res-cutpoints} to~\ref{sec:res-selection} report the three ablations of the
fine-tuning track, Section~\ref{sec:res-softqwk} the two-phase Soft-QWK method,
Section~\ref{sec:res-newloss} the three losses added from readability research in other
languages, and Section~\ref{sec:res-frozen} the frozen-backbone track with its own ablation. Section~\ref{sec:res-context}
places the numbers beside the published systems and beside human agreement, and
Section~\ref{sec:res-limits} states the limitations.

Every number in this chapter is read from the result files produced by the training runs. The
values quoted in the text and the values in the tables are generated from those files by a
script, so the text and the tables cannot disagree.

\section{Integrity of the Runs}
\label{sec:res-integrity}

The three guards declared in Section~\ref{sec:meth-impl} passed identically in all eleven runs,
and the result files record the evidence.

\begin{itemize}
\item Exactly one accelerator was visible in every run, an NVIDIA \gpuName, and the realised
training batch was 64 sentences in every run, including both phases of the two-phase method.
\item Every epoch of every run performed \stepsPerEpoch\ optimiser steps, which is the number
implied by 54{,}845 training sentences at an effective batch of 64.
\item The reloaded selected checkpoint reproduced the development QWK, accuracy and distance
logged during training at that epoch, to within $10^{-9}$, in every run. The numbers reported
below therefore belong to the epoch the selection rule chose.
\item All eleven runs used the same corpus revision (\datasetSha) and the same encoder revision
(\modelSha), with Python \pythonVersion, PyTorch \torchVersion\ and Transformers \tfVersion,
and none of them was a reduced smoke run.
\end{itemize}

The eleven runs together took \totalMinutes\ minutes of accelerator time, between \minAll\ and
\maxAll\ minutes each.

\section{Reproduction of the Benchmark}
\label{sec:res-repro}

Eight of the nine methods in the comparable group are loss functions that the BAREC benchmark
also reports, on the same corpus, the same encoder and the same input representation. Comparing
them is therefore a test of the pipeline rather than a contribution, and it has to be passed
before anything else in this chapter can be believed.
Table~\ref{tab:res-repro} gives the comparison on the development set.

% generated by gen_numbers.py
\begin{table}[H]
\centering
\footnotesize
\caption{Reproduction of the eight loss functions of BAREC Table~6 on the development set. Paper values are from Elmadani et al.~\cite{elmadani2025barec}; the gap is ours minus theirs.}
\label{tab:res-repro}
\begin{tabular}{lrrrrrrrr}
\toprule
& \multicolumn{2}{c}{\textbf{QWK}} & \multicolumn{2}{c}{\textbf{Acc19}} & \multicolumn{2}{c}{\textbf{$\pm$1 Acc19}} & \multicolumn{2}{c}{\textbf{Dist}} \\
\cmidrule(lr){2-3}\cmidrule(lr){4-5}\cmidrule(lr){6-7}\cmidrule(lr){8-9}
\textbf{Loss} & ours & paper & ours & paper & ours & paper & ours & paper \\
\midrule
CE & 79.81 & 80.0 & 54.84 & 56.6 & 70.47 & 69.9 & 1.163 & 1.14 \\
Regression (MSE) & 82.88 & 84.0 & 35.54 & 43.1 & 71.15 & 73.1 & 1.212 & 1.13 \\
OLL ($\alpha=1.5$) & 82.91 & 82.8 & 46.21 & 47.3 & 71.15 & 71.1 & 1.130 & 1.13 \\
EMD & 81.29 & 81.2 & 54.27 & 55.3 & 70.10 & 70.3 & 1.125 & 1.11 \\
SOFT ($\beta=2$) & 81.41 & 80.0 & 57.07 & 55.8 & 71.66 & 69.8 & 1.088 & 1.15 \\
SOFT ($\beta=3$) & 79.72 & 80.1 & 55.40 & 56.4 & 69.44 & 69.9 & 1.177 & 1.14 \\
SOFT ($\beta=4$) & 80.05 & 79.6 & 55.62 & 56.4 & 70.79 & 69.9 & 1.144 & 1.15 \\
SOFT ($\beta=5$) & 80.63 & 79.3 & 55.85 & 56.2 & 70.86 & 69.5 & 1.134 & 1.17 \\
\bottomrule
\end{tabular}
\end{table}


Seven of the eight losses land within 1.5 QWK points of the published value, and the mean
absolute difference over the eight is \gapQwkMean\ points. Cross-entropy reproduces to
\gapQwkCE\ points, EMD to \gapQwkEMD\ and OLL to \gapQwkOLL, which is as close as a single run
of a fine-tuning pipeline can reasonably be expected to land. The soft-label family sits above
the published values, most clearly at $\beta=2$ (\gapQwkSoftTwo) and $\beta=5$
(\gapQwkSoftFive), and its ordering across $\beta$ is not the published one. The four values
span 1.69 QWK points in our run while the published four span 0.8, so this is the family in
which a single run is least informative.

The one substantial difference is the regression method, at \gapQwkMSE\ QWK points and
\gapAccMSE\ points of exact accuracy. Table~\ref{tab:res-epochs} explains most of it.

% generated by gen_numbers.py
\begin{table}[H]
\centering
\footnotesize
\caption{Development QWK after each epoch. The selected epoch, that is the epoch with the lowest development loss, is starred. The two-phase method is shown as its three mean squared error epochs followed by its three Soft-QWK epochs.}
\label{tab:res-epochs}
\begin{tabular}{lrrrrrr}
\toprule
\textbf{Method} & \textbf{ep.\ 1} & \textbf{2} & \textbf{3} & \textbf{4} & \textbf{5} & \textbf{6} \\
\midrule
CE & 79.12 & \textbf{79.81}$^{\star}$ & 81.12 & 80.40 & 80.83 & 80.45 \\
Regression (MSE) & 81.90 & \textbf{82.88}$^{\star}$ & 83.51 & 83.55 & 82.83 & 83.08 \\
OLL ($\alpha=1.5$) & 81.71 & 82.57 & \textbf{82.91}$^{\star}$ & 82.33 & 82.41 & 82.43 \\
EMD & 79.63 & \textbf{81.29}$^{\star}$ & 81.07 & 81.08 & 80.80 & 81.28 \\
SOFT ($\beta=2$) & 79.96 & \textbf{81.41}$^{\star}$ & 81.58 & 81.57 & 80.69 & 80.70 \\
SOFT ($\beta=3$) & \textbf{79.72}$^{\star}$ & 80.69 & 80.89 & 80.71 & 80.15 & 80.46 \\
SOFT ($\beta=4$) & 78.74 & \textbf{80.05}$^{\star}$ & 81.09 & 80.77 & 80.87 & 80.45 \\
SOFT ($\beta=5$) & 79.73 & \textbf{80.63}$^{\star}$ & 81.18 & 81.12 & 80.86 & 80.63 \\
CORAL & 81.23 & \textbf{82.90}$^{\star}$ & 83.11 & 82.32 & 82.07 & 82.09 \\
CORAL, head lr $10^{-3}$ & 81.93 & 83.41 & \textbf{83.86}$^{\star}$ & 82.79 & 82.41 & 82.42 \\
\midrule
MSE warm-up (phase 1) & 81.24 & 82.54 & 83.10 & \multicolumn{3}{c}{--} \\
Soft-QWK (phase 2) & \multicolumn{3}{c}{--} & 82.13 & \textbf{83.24}$^{\star}$ & 82.57 \\
\bottomrule
\end{tabular}
\end{table}


The development QWK of the regression method keeps rising until epoch four, but its
development \emph{loss} is lowest at epoch \selMSE, and the selection rule takes that epoch. At
epoch four the method reaches \qwkseldevQwkMSE\ development QWK and \qwkseldevAccMSE\ exact
accuracy, against the published \paperQwkMSE\ and \paperAccMSE; on the
test set it reaches \qwkseltestQwkMSE\ QWK and \qwkseltestAccMSE\ exact accuracy against the
published 84.4 and 41.4. In other words the implementation does reach the published numbers,
and the gap in Table~\ref{tab:res-repro} is the price of one checkpoint decision, taken by a
rule that both the benchmark and this thesis apply. The mean squared error loss is only loosely
coupled to QWK, so the epoch that minimises it need not be the epoch that maximises the metric;
the same effect appears, smaller, for several other methods and is quantified in
Section~\ref{sec:res-selection}.

Table~\ref{tab:res-collapsed} repeats the comparison at the coarser granularities that the
benchmark reports for three of the losses.

% generated by gen_numbers.py
\begin{table}[H]
\centering
\footnotesize
\caption{Collapsed granularities for the three losses the BAREC paper reports at 7, 5 and 3 levels (Tables~7 and~8 of Elmadani et al.~\cite{elmadani2025barec}).}
\label{tab:res-collapsed}
\begin{tabular}{llrrrrrr}
\toprule
& & \multicolumn{2}{c}{\textbf{Acc7}} & \multicolumn{2}{c}{\textbf{Acc5}} & \multicolumn{2}{c}{\textbf{Acc3}} \\
\cmidrule(lr){3-4}\cmidrule(lr){5-6}\cmidrule(lr){7-8}
\textbf{Split} & \textbf{Loss} & ours & paper & ours & paper & ours & paper \\
\midrule
dev & CE & 65.16 & 65.9 & 69.96 & 70.3 & 76.76 & 76.5 \\
dev & EMD & 65.01 & 65.2 & 70.10 & 70.0 & 76.63 & 76.4 \\
dev & Regression (MSE) & 57.02 & 61.1 & 64.43 & 67.8 & 73.34 & 75.9 \\
\midrule
test & CE & 64.36 & 65.1 & 68.71 & 69.4 & 74.81 & 75.2 \\
test & EMD & 64.00 & 64.9 & 68.97 & 69.0 & 74.87 & 75.2 \\
test & Regression (MSE) & 56.24 & 59.4 & 63.05 & 65.3 & 71.31 & 72.8 \\
\bottomrule
\end{tabular}
\end{table}


Cross-entropy and EMD agree with the published collapsed accuracies to within about one point on
both splits. The regression method is again lower, by three to four points, which is consistent
with the same checkpoint effect: collapsed accuracy is an exact-match measure, and the epoch
chosen for the regression method is the one at which its predictions are still concentrated
near the centre of the scale.

The conclusion of this section is that the pipeline reproduces the benchmark. That licence is
what allows the next sections to add CORAL and Soft-QWK to the same table and to read the
differences as differences between methods.

\section{Main Results}
\label{sec:res-main}

Tables~\ref{tab:res-main-dev} and~\ref{tab:res-main-test} report every metric for every run on
the development and test splits.

% generated by gen_numbers.py
\begin{table}[H]
\centering
\footnotesize
\caption{All methods on the BAREC development set. AraBERTv2 with D3Tok input, one seed.}
\label{tab:res-main-dev}
\begin{adjustbox}{max width=\textwidth}
\begin{tabular}{lrrrrrrrrrrr}
\toprule
\textbf{Method} & \textbf{QWK} & \textbf{Acc19} & \textbf{$\pm$1} & \textbf{Dist} & \textbf{Acc7} & \textbf{Acc5} & \textbf{Acc3} & \textbf{far$\geq$3} & \textbf{far$\geq$5} & \textbf{mF1} & \textbf{Ep.} \\
\midrule
\multicolumn{12}{l}{\textit{Comparable group: identical protocol, one head and one loss changed}} \\
\midrule
CE & 79.81 & 54.84 & 70.47 & 1.163 & 65.16 & 69.96 & 76.76 & 15.60 & 6.69 & 45.45 & 2 \\
Regression (MSE) & 82.88 & 35.54 & 71.15 & 1.212 & 57.02 & 64.43 & 73.34 & 13.43 & 3.67 & 22.13 & 2 \\
OLL ($\alpha=1.5$) & 82.91 & 46.21 & 71.15 & 1.130 & 61.55 & 67.62 & 76.33 & 13.45 & 4.24 & 27.64 & 3 \\
EMD & 81.29 & 54.27 & 70.10 & 1.125 & 65.01 & 70.10 & 76.63 & 14.12 & 5.84 & 36.60 & 2 \\
SOFT ($\beta=2$) & 81.41 & 57.07 & 71.66 & 1.088 & 66.95 & 71.23 & 77.44 & 14.23 & 5.84 & 46.88 & 2 \\
SOFT ($\beta=3$) & 79.72 & 55.40 & 69.44 & 1.177 & 65.20 & 69.59 & 76.01 & 15.25 & 6.92 & 40.05 & 1 \\
SOFT ($\beta=4$) & 80.05 & 55.62 & 70.79 & 1.144 & 65.94 & 70.21 & 76.84 & 15.17 & 6.31 & 43.88 & 2 \\
SOFT ($\beta=5$) & 80.63 & 55.85 & 70.86 & 1.134 & 66.21 & 70.26 & 77.02 & 15.23 & 6.32 & 44.79 & 2 \\
CORAL & 82.90 & 34.15 & 66.84 & 1.289 & 55.40 & 63.60 & 73.00 & 12.86 & 4.32 & 24.56 & 2 \\
\midrule
\multicolumn{12}{l}{\textit{Two-phase group: MSE warm-up then Soft-QWK, so the schedule differs}} \\
\midrule
MSE $\rightarrow$ Soft-QWK & 83.24 & 35.23 & 72.33 & 1.251 & 58.29 & 65.76 & 75.31 & 13.83 & 4.73 & 28.03 & 2 \\
\midrule
\multicolumn{12}{l}{\textit{Ablation: not part of the comparable group}} \\
\midrule
CORAL, head lr $10^{-3}$ & 83.86 & 42.56 & 69.37 & 1.167 & 59.02 & 66.25 & 75.25 & 12.45 & 4.13 & 28.88 & 3 \\
\bottomrule
\end{tabular}
\end{adjustbox}

\vspace{0.1cm}
{\footnotesize QWK, the accuracies, far$\geq$3, far$\geq$5 and macro-F1 are percentages; Dist is in levels; Ep.\ is the selected epoch. far$\geq$3 and far$\geq$5 are this thesis's own metrics.}
\end{table}


% generated by gen_numbers.py
\begin{table}[H]
\centering
\footnotesize
\caption{All methods on the BAREC test set. Same checkpoints as Table~\ref{tab:res-main-dev}.}
\label{tab:res-main-test}
\begin{adjustbox}{max width=\textwidth}
\begin{tabular}{lrrrrrrrrrrr}
\toprule
\textbf{Method} & \textbf{QWK} & \textbf{Acc19} & \textbf{$\pm$1} & \textbf{Dist} & \textbf{Acc7} & \textbf{Acc5} & \textbf{Acc3} & \textbf{far$\geq$3} & \textbf{far$\geq$5} & \textbf{mF1} & \textbf{Ep.} \\
\midrule
\multicolumn{12}{l}{\textit{Comparable group: identical protocol, one head and one loss changed}} \\
\midrule
CE & 80.48 & 54.72 & 70.29 & 1.128 & 64.36 & 68.71 & 74.81 & 15.06 & 6.30 & 46.86 & 2 \\
Regression (MSE) & 83.64 & 34.20 & 71.70 & 1.186 & 56.24 & 63.05 & 71.31 & 13.41 & 2.88 & 23.61 & 2 \\
OLL ($\alpha=1.5$) & 84.08 & 46.28 & 72.08 & 1.074 & 60.69 & 66.09 & 73.80 & 12.52 & 3.60 & 29.01 & 3 \\
EMD & 82.70 & 53.58 & 69.39 & 1.080 & 64.00 & 68.97 & 74.87 & 13.33 & 5.23 & 38.18 & 2 \\
SOFT ($\beta=2$) & 81.86 & 55.96 & 71.18 & 1.070 & 65.43 & 69.60 & 75.60 & 14.00 & 5.56 & 47.19 & 2 \\
SOFT ($\beta=3$) & 80.31 & 54.64 & 69.00 & 1.142 & 63.88 & 68.24 & 74.70 & 14.69 & 6.49 & 42.14 & 1 \\
SOFT ($\beta=4$) & 80.94 & 55.61 & 70.93 & 1.095 & 65.18 & 69.39 & 75.20 & 14.40 & 5.82 & 46.04 & 2 \\
SOFT ($\beta=5$) & 81.37 & 56.03 & 70.93 & 1.086 & 65.41 & 69.60 & 75.50 & 14.52 & 5.87 & 46.37 & 2 \\
CORAL & 84.33 & 34.15 & 66.70 & 1.232 & 54.09 & 60.98 & 69.21 & 12.06 & 3.29 & 26.38 & 2 \\
\midrule
\multicolumn{12}{l}{\textit{Two-phase group: MSE warm-up then Soft-QWK, so the schedule differs}} \\
\midrule
MSE $\rightarrow$ Soft-QWK & 84.31 & 35.96 & 72.21 & 1.196 & 57.33 & 64.36 & 72.80 & 13.07 & 4.19 & 28.96 & 2 \\
\midrule
\multicolumn{12}{l}{\textit{Ablation: not part of the comparable group}} \\
\midrule
CORAL, head lr $10^{-3}$ & 84.58 & 42.15 & 70.37 & 1.121 & 58.39 & 64.59 & 73.03 & 11.57 & 3.49 & 28.70 & 3 \\
\bottomrule
\end{tabular}
\end{adjustbox}

\vspace{0.1cm}
{\footnotesize QWK, the accuracies, far$\geq$3, far$\geq$5 and macro-F1 are percentages; Dist is in levels; Ep.\ is the selected epoch. far$\geq$3 and far$\geq$5 are this thesis's own metrics.}
\end{table}


\subsection{CORAL and the QWK-Oriented Objective}

The question in the research brief is whether an ordinal head and a metric-aligned loss beat
distance-aware losses when the dataset, the protocol and the metric are held fixed. The answer
is that they are at the top of the table, and that they are not alone there.

On the development set the comparable group is led by OLL at \devQwkOLL, CORAL at
\devQwkCoral\ and the regression method at \devQwkMSE. The distance between these three is
0.03 QWK points. The two-phase Soft-QWK method reaches \devQwkSoftqwk, which is
\softqwkPhaseTwoGain\ above its own warm-up checkpoint and the highest development value
outside the ablation, but it is a different training schedule and is therefore reported
separately. On the
test set the order of the leaders changes: CORAL reaches \testQwkCoral, Soft-QWK
\testQwkSoftqwk, OLL \testQwkOLL\ and the regression method \testQwkMSE.

Two things follow, and they should be stated with the same care.

First, ordinal-aware objectives clearly beat the ordinal-blind baseline. Cross-entropy is
\devQwkCE\ on development and \testQwkCE\ on test, which together with SOFT at $\beta=3$
forms the bottom of the comparable group. Every other method is above it, by up to three QWK
points on development and up to four on test; the single soft-label configuration that does not
beat it is $\beta=3$, at \devQwkSoftThree\ and \testQwkSoftThree. This is the same
conclusion the benchmark draws, and it is reproduced here under a protocol in which nothing but
the head and the loss changed.

Second, the differences \emph{among} the ordinal-aware objectives are small, and the ranking is
not stable between splits: the leader on development is fourth on test, and CORAL rises from
\devQwkCoral\ on development to \testQwkCoral\ on test while the regression method rises by
less than a point. With one seed per method, differences of a few tenths of a QWK point cannot be attributed
to the mechanism. This thesis therefore does not name a single best loss. What the evidence
supports is a group of four objectives, namely the CORAL head, the metric-aligned two-phase
schedule, distance-weighted ordinal log loss and plain regression, that are equivalent within
the resolution of this experiment and that all clearly beat cross-entropy and the soft-label
family on QWK.

\subsection{Where the Groups Differ}

Reading across Table~\ref{tab:res-main-dev} shows that the four leaders on QWK are not the same
methods that lead on the other metrics, and that the split is systematic. The two methods with
a continuous or cumulative output, that is the regression head and the CORAL head, have the
highest QWK and the lowest exact accuracy: \devAccMSE\ and \devAccCoral\ respectively, against
\devAccSoftTwo\ for the best soft-label configuration. Their macro-F1 values, \devMfMSE\ and
\devMfCoral\ against \devMfSoftTwo, say the same thing in a different way: they concentrate
their predictions on a narrower band of levels and rarely commit to a rare level. The
classification head with a distance-aware loss does the opposite: it spreads predictions across
the whole scale, gets more of them exactly right, and pays for the attempts with a longer tail
of errors. Section~\ref{sec:res-tradeoff} quantifies this trade-off and
Section~\ref{sec:res-distance} looks at the tail itself.

\section{Exact Match against Ordinal Agreement}
\label{sec:res-tradeoff}

The trade-off is visible as an almost exact inversion between two columns of the same table. Among
the ten runs of the comparable and two-phase groups, the four highest development QWK values
belong to Soft-QWK, OLL, CORAL and regression, and those same four hold the four lowest exact
accuracies: \devAccCoral\ for CORAL, \devAccSoftqwk\ for Soft-QWK, \devAccMSE\ for regression
and \devAccOLL\ for OLL. Conversely the six methods with the highest exact accuracy, led by
SOFT at $\beta=2$ with \devAccSoftTwo, occupy the lower half of the QWK column.

The mechanism is not mysterious. A regression or cumulative model is rewarded for placing its
prediction close to the gold level and is never rewarded for landing on it exactly, so it
learns a smooth, cautious mapping and keeps away from the ends of the scale. A softmax model
with a distance-aware loss is still trained to put mass on the gold level, so it keeps
committing to specific levels, wins exact matches, and occasionally commits badly.

This has a practical consequence that the single headline metric hides. If the intended use is
to place a text in a narrow band, for example to choose which grade a text suits, then QWK and
the distance metrics are the right criteria and the regression or CORAL family is the right
choice. If the intended use needs the exact level, for instance to label material by a
curriculum level without human review, then exact accuracy matters and the soft-label family is
nearly twenty points better at it than the QWK leaders, at a cost of two to three QWK points.
Reporting only QWK, as both the benchmark and the shared task do, obscures a difference this
large.

\section{Error Distance: Near Misses and Far Misses}
\label{sec:res-distance}

The research brief asks not only which method scores best but which produces the smallest error
distance, and specifically whether errors are near misses rather than far misses. Four columns
of Tables~\ref{tab:res-main-dev} and~\ref{tab:res-main-test} answer that: adjacent accuracy,
mean distance, and the two far-miss rates introduced in this thesis.

The first observation is that the best method on QWK is not the best method on distance. On the
development set the lowest mean distance in the study is \devDistSoftTwo\ levels, from SOFT at
$\beta=2$, whose QWK is \devQwkSoftTwo, two points below the leaders. The QWK leaders have
\emph{worse} mean distance: \devDistCoral\ for CORAL and \devDistSoftqwk\ for Soft-QWK, the two
worst values in the whole table. The same pattern holds on test, where the lowest mean distance
is \testDistSoftTwo\ and CORAL is at \testDistCoral.

This looks contradictory only until the two tails are separated. Mean distance averages all
errors, so it is dominated by the many small ones; QWK weights an error by its square, so it is
dominated by the few large ones. The far-miss rates show which methods control the large
errors. On the development set the lowest rate of errors of three levels or more is
\devFarThreeCoral\ per cent, from CORAL, followed by the regression method at
\devFarThreeMSE\ and OLL at \devFarThreeOLL; cross-entropy is worst at \devFarThreeCE. For
errors of five levels or more the order is the regression method at \devFarFiveMSE\ per cent,
OLL at \devFarFiveOLL\ and CORAL at \devFarFiveCoral, against \devFarFiveCE\ for cross-entropy.
On the test set CORAL again has the fewest errors of three levels or more, \testFarThreeCoral\
per cent, and the regression method the fewest of five levels or more, \testFarFiveMSE.

So the ordinal mechanisms do exactly what they were designed to do, and the effect is specific:
cross-entropy makes roughly one error of five levels or more in every fifteen sentences, while
the regression head makes one in every twenty-seven on the development set. What the ordinal
mechanisms do not do is reduce the number of small errors; the soft-label family, which keeps a
cross-entropy-like target, remains better at that.

Adjacent accuracy behaves like a third, intermediate criterion. The highest development value is
\devPmSoftqwk\ per cent for Soft-QWK, followed by \devPmSoftTwo\ for SOFT at $\beta=2$ and
\devPmMSE\ for regression, while CORAL is last at \devPmCoral. CORAL therefore achieves its
high QWK without being good at landing within one level: it avoids the catastrophic errors that
QWK punishes hardest, but it also misses the adjacent band more often than any other method.
For an application in which a one-level error is acceptable and a three-level error is not,
CORAL and the regression method are the appropriate choices; if a one-level error already has a
cost, the soft-label family is preferable.

Taken together, the three views of error distance answer the brief's question in a way a single
metric cannot: the methods that maximise QWK are the methods that minimise far misses, not the
methods that minimise average error.

\section{CORAL Cut-Points and the Travel Budget}
\label{sec:res-cutpoints}

Section~\ref{sec:meth-coral-init} argued that CORAL's cut-points cannot travel far under the
BAREC protocol, and initialised them from the training label frequencies for that reason. Two
measurements test the argument.

The first is the distance the cut-points actually moved. Over the whole run, no cut-point moved
further than \cutMoveFineCoral\ from its initial value, against a fitted spread of
\cutSpreadCoral\ between the highest and the lowest cut-point. The budget computed in
Section~\ref{sec:meth-coral-init} was \lrSumEncoder, and the observed movement is well inside
it. The trained cut-points run from \cutTopCoral\ down to \cutBottomCoral, which is to say they
are, to two decimal places, the values that Equation~\ref{eq:coral-init} wrote into them before
training started. Under this protocol the cut-points are not learned; they are supplied. A
zero-initialised head would have had to move each of them by up to six units on a budget of
\lrSumEncoder, which the arithmetic forbids. Table~\ref{tab:res-cutpoints} also confirms that
the 18 cut-points stayed correctly ordered after every epoch, so the rank consistency that
motivates the CORAL head held throughout training and not only at the optimum.

% generated by gen_numbers.py
\begin{table}[H]
\centering
\footnotesize
\caption{CORAL cut-point spread $b_0-b_{17}$ after each epoch, and whether the 18 cut-points were still ordered. Measured on the saved model at the end of every epoch; the largest distance any cut-point ever travelled from its initial value is \cutMoveFineCoral\ with the shared learning rate and \cutMoveCoralHeadlr\ in the head learning rate ablation.}
\label{tab:res-cutpoints}
\begin{tabular}{lrrrrrrr}
\toprule
\textbf{Method} & \textbf{ep.\ 1} & \textbf{2} & \textbf{3} & \textbf{4} & \textbf{5} & \textbf{6} & \textbf{ordered} \\
\midrule
CORAL & 11.62 & 11.68 & 11.73 & 11.76 & 11.78 & 11.79 & yes \\
CORAL, head lr $10^{-3}$ & 12.65 & 13.62 & 14.43 & 15.01 & 15.38 & 15.51 & yes \\
\bottomrule
\end{tabular}
\end{table}


The second measurement is the ablation that raises the head learning rate to $10^{-3}$ and
leaves the encoder at $5\times10^{-5}$, multiplying the cut-point budget by twenty to
\lrSumHeadlr. The cut-points then move, by up to \cutMoveCoralHeadlr, and their spread widens
from \cutSpreadCoral\ to \cutSpreadCoralHeadlr. Development QWK rises by
\coralHeadlrDevGain\ points to \devQwkCoralHeadlr\ and test QWK by \coralHeadlrTestGain\ points
to \testQwkCoralHeadlr, the highest test value recorded in this work. Exact accuracy rises as
well, from \devAccCoral\ to \devAccCoralHeadlr, and the far-miss rate falls from
\devFarThreeCoral\ to \devFarThreeCoralHeadlr.

Both measurements are consistent with the argument, and the ablation adds something the
argument did not predict: the cut-points are not merely hard to move, they are worth moving.
Since the ablation changes an optimiser setting rather than the head, it is reported outside the
comparable group; it is also a single run, so the size of its gain should be treated as an
indication and not as a measured effect. The honest summary is that the initialisation of
Equation~\ref{eq:coral-init} makes a CORAL head usable under a protocol whose learning rate was
chosen for an encoder, and that a head-specific learning rate appears to be the better remedy.

\section{Decision Threshold for CORAL}
\label{sec:res-threshold}

The ZAI system decodes its CORAL head with a threshold tuned on the development set rather than
with $0.5$ \cite{nazzal2025zai}. Repeating that on the saved CORAL model selects
$t=\zaiTCoral$, barely away from the default, and moves test QWK by \zaiTestGainCoral\ points,
from \testQwkCoral\ to \zaiTestQwkCoral. In the head learning rate ablation the grid selects
$t=\zaiTCoralHeadlr$ and nothing changes.

The threshold is therefore not where the ZAI system's advantage comes from, at least not with
cut-points initialised as they are here. This is what Equation~\ref{eq:coral-init} predicts: an
initialisation that already encodes the cumulative label frequencies places the decision
boundaries where the data puts them, so there is little left for a global threshold shift to
correct.

\section{The Checkpoint Selection Rule}
\label{sec:res-selection}

Both shared-task systems select their checkpoint on development QWK
\cite{nazzal2025zai,trigui2025stbw}, while the benchmark, and therefore this thesis, selects on
development loss. Table~\ref{tab:res-selection} evaluates both rules for every method. Only the
test columns can be read as results, since a checkpoint chosen on development QWK has an
optimistic development QWK by construction.

% generated by gen_numbers.py
\begin{table}[H]
\centering
\footnotesize
\caption{Checkpoint selection ablation. The left block is BAREC's rule (lowest development loss), the right block the rule used by the two shared-task systems (highest development QWK). Development QWK under the right rule is optimistic by construction, so only the test column should be read as a result.}
\label{tab:res-selection}
\begin{adjustbox}{max width=\textwidth}
\begin{tabular}{lrrrrrrr}
\toprule
& \multicolumn{3}{c}{\textbf{lowest dev loss}} & \multicolumn{3}{c}{\textbf{highest dev QWK}} & \\
\cmidrule(lr){2-4}\cmidrule(lr){5-7}
\textbf{Method} & epoch & dev QWK & test QWK & epoch & dev QWK & test QWK & \textbf{$\Delta$ test} \\
\midrule
CE & 2 & 79.81 & 80.48 & 3 & 81.12 & 81.80 & +1.32 \\
Regression (MSE) & 2 & 82.88 & 83.64 & 4 & 83.55 & 84.32 & +0.67 \\
OLL ($\alpha=1.5$) & 3 & 82.91 & 84.08 & 3 & 82.91 & 84.08 & +0.00 \\
EMD & 2 & 81.29 & 82.70 & 2 & 81.29 & 82.70 & +0.00 \\
SOFT ($\beta=2$) & 2 & 81.41 & 81.86 & 3 & 81.58 & 82.38 & +0.52 \\
SOFT ($\beta=3$) & 1 & 79.72 & 80.31 & 3 & 80.89 & 82.31 & +2.00 \\
SOFT ($\beta=4$) & 2 & 80.05 & 80.94 & 3 & 81.09 & 82.61 & +1.67 \\
SOFT ($\beta=5$) & 2 & 80.63 & 81.37 & 3 & 81.18 & 82.19 & +0.82 \\
CORAL & 2 & 82.90 & 84.33 & 3 & 83.11 & 84.50 & +0.17 \\
MSE $\rightarrow$ Soft-QWK & 2 & 83.24 & 84.31 & 2 & 83.24 & 84.31 & +0.00 \\
CORAL, head lr $10^{-3}$ & 3 & 83.86 & 84.58 & 3 & 83.86 & 84.58 & +0.00 \\
\bottomrule
\end{tabular}
\end{adjustbox}
\end{table}


The rule is worth up to \qwkselTestGainMax\ QWK points on the test set. It changes nothing for
the four runs whose loss and QWK already peak at the same epoch, namely OLL, EMD, the two-phase
method and the head learning rate ablation. It is worth \qwkselTestGainCE\ points for cross-entropy,
\qwkselTestGainSoftThree\ for SOFT at $\beta=3$ and \qwkselTestGainSoftFour\ for SOFT at
$\beta=4$, and only \qwkselTestGainCoral\ for CORAL and \qwkselTestGainMSE\ for the regression
method.

Two conclusions follow. First, part of the difference between published systems that use
different selection rules is the rule and not the model, and a comparison that does not hold the
rule fixed is not a comparison of methods. Second, the effect is not uniform: it is largest for
the methods whose loss is least aligned with the metric, and smallest for CORAL and regression,
whose losses already reward proximity. The gaps between the methods in
Table~\ref{tab:res-main-test} shrink under QWK selection, because the methods that gain most
are the ones that were behind.

\section{The Two-Phase Soft-QWK Method}
\label{sec:res-softqwk}

The two-phase method was trained as mean squared error for three epochs and then as Soft-QWK
for three more, on the same six-epoch budget as every other run. The warm-up checkpoint, chosen
by the same rule inside phase one, reaches \phaseOneQwk\ development QWK with \phaseOneAcc\
exact accuracy and a mean distance of \phaseOneDist. After the second phase the same model
reaches \devQwkSoftqwk\ development QWK, that is \softqwkPhaseTwoGain\ points, with adjacent
accuracy rising to \devPmSoftqwk, the highest in the study, and the mean distance rising to
\devDistSoftqwk, one of the two worst. On the test set the method reaches \testQwkSoftqwk.

The size and the direction of that gain match what the STBW authors report for their own system,
which improved from 83.9 to 84.88 QWK when the second phase was added
\cite{trigui2025stbw}. Their absolute numbers are not comparable with ours, for two reasons:
they evaluate on the shared-task splits, and their input carries corpus metadata and
surface statistics in addition to the segmented text, while the input here is the segmented
text alone.

Table~\ref{tab:res-epochs} shows the phase-two trajectory, which rises from 82.13 to
\devQwkSoftqwk\ and then falls to 82.57. A batch-level objective that can be lowered by
predicting near the centre of the scale is expected to drift in exactly that way, and the
metrics of the selected epoch show the drift beginning: adjacent accuracy is the best in the
study while the mean distance and the exact accuracy are among the worst, which is the profile
of a model that has learned to stay close to the middle. The warm-up is therefore not a
formality, and the fact that this thesis reproduces the benefit of the second phase under a
protocol that gives it no extra epochs is the substantive result here.

\section{Loss Functions from Readability Research in Other Languages}
\label{sec:res-newloss}

This section reports the three losses of Section~\ref{sec:meth-newlosses}: CORN, and the two
pairwise ranking objectives added to the regression loss, margin ranking (RANK) and pairwise
logistic (PLR). They were trained after the main comparison, with the same script and protocol,
and the same integrity checks passed in all of them: one accelerator, a realised batch of
64 sentences, \stepsPerEpoch\ optimiser steps in every epoch, and a reloaded checkpoint that
reproduces the logged development metrics exactly. Tables~\ref{tab:res-newloss-dev}
and~\ref{tab:res-newloss-test} give their results beside their matched controls, and
Table~\ref{tab:res-newloss-delta} gives each difference to its control.

% generated by gen_newloss_numbers.py
\begin{table}[H]
\centering
\footnotesize
\caption{Loss functions from readability research in other languages, on the BAREC development set, beside their matched controls from Table~\ref{tab:res-main-dev}. Same pipeline, one head and one loss changed.}
\label{tab:res-newloss-dev}
\begin{adjustbox}{max width=\textwidth}
\begin{tabular}{lrrrrrrrrrrr}
\toprule
\textbf{Method} & \textbf{QWK} & \textbf{Acc19} & \textbf{$\pm$1} & \textbf{Dist} & \textbf{Acc7} & \textbf{Acc5} & \textbf{Acc3} & \textbf{far$\geq$3} & \textbf{far$\geq$5} & \textbf{mF1} & \textbf{Ep.} \\
\midrule
\multicolumn{12}{l}{\textit{Controls from the comparable group}} \\
\midrule
Regression (MSE) & 82.88 & 35.54 & 71.15 & 1.212 & 57.02 & 64.43 & 73.34 & 13.43 & 3.67 & 22.13 & 2 \\
CORAL & 82.90 & 34.15 & 66.84 & 1.289 & 55.40 & 63.60 & 73.00 & 12.86 & 4.32 & 24.56 & 2 \\
\midrule
\multicolumn{12}{l}{\textit{Loss functions added in this thesis}} \\
\midrule
CORN & 81.92 & 54.47 & 71.46 & 1.096 & 65.39 & 70.05 & 76.79 & 13.83 & 5.34 & 42.57 & 3 \\
MSE + margin ranking (RANK) & 82.67 & 37.81 & 71.08 & 1.209 & 58.28 & 65.98 & 74.97 & 13.53 & 4.06 & 23.03 & 3 \\
MSE + pairwise logistic (PLR) & 82.61 & 30.70 & 69.77 & 1.269 & 55.50 & 63.52 & 73.89 & 14.30 & 3.34 & 18.21 & 2 \\
\bottomrule
\end{tabular}
\end{adjustbox}

\vspace{0.1cm}
{\footnotesize QWK, the accuracies, far$\geq$3, far$\geq$5 and macro-F1 are percentages; Dist is in levels; Ep.\ is the selected epoch. far$\geq$3 and far$\geq$5 are this thesis's own metrics.}
\end{table}


% generated by gen_newloss_numbers.py
\begin{table}[H]
\centering
\footnotesize
\caption{Loss functions from readability research in other languages, on the BAREC test set, beside their matched controls from Table~\ref{tab:res-main-test}. Same pipeline, one head and one loss changed.}
\label{tab:res-newloss-test}
\begin{adjustbox}{max width=\textwidth}
\begin{tabular}{lrrrrrrrrrrr}
\toprule
\textbf{Method} & \textbf{QWK} & \textbf{Acc19} & \textbf{$\pm$1} & \textbf{Dist} & \textbf{Acc7} & \textbf{Acc5} & \textbf{Acc3} & \textbf{far$\geq$3} & \textbf{far$\geq$5} & \textbf{mF1} & \textbf{Ep.} \\
\midrule
\multicolumn{12}{l}{\textit{Controls from the comparable group}} \\
\midrule
Regression (MSE) & 83.64 & 34.20 & 71.70 & 1.186 & 56.24 & 63.05 & 71.31 & 13.41 & 2.88 & 23.61 & 2 \\
CORAL & 84.33 & 34.15 & 66.70 & 1.232 & 54.09 & 60.98 & 69.21 & 12.06 & 3.29 & 26.38 & 2 \\
\midrule
\multicolumn{12}{l}{\textit{Loss functions added in this thesis}} \\
\midrule
CORN & 82.57 & 54.39 & 71.49 & 1.058 & 64.69 & 69.02 & 75.20 & 13.34 & 4.87 & 43.97 & 3 \\
MSE + margin ranking (RANK) & 84.06 & 36.91 & 71.96 & 1.158 & 57.53 & 64.36 & 72.43 & 12.41 & 3.31 & 23.34 & 3 \\
MSE + pairwise logistic (PLR) & 83.68 & 30.19 & 70.74 & 1.219 & 56.03 & 63.19 & 71.89 & 13.08 & 2.54 & 18.63 & 2 \\
\bottomrule
\end{tabular}
\end{adjustbox}

\vspace{0.1cm}
{\footnotesize QWK, the accuracies, far$\geq$3, far$\geq$5 and macro-F1 are percentages; Dist is in levels; Ep.\ is the selected epoch. far$\geq$3 and far$\geq$5 are this thesis's own metrics.}
\end{table}


% generated by gen_newloss_numbers.py
\begin{table}[H]
\centering
\footnotesize
\caption{Each added loss minus its matched control, in percentage points. A positive QWK or Acc19 difference is better; a negative far-miss difference is better. One seed per method.}
\label{tab:res-newloss-delta}
\begin{adjustbox}{max width=\textwidth}
\begin{tabular}{llrrrrrrrr}
\toprule
& & \multicolumn{4}{c}{\textbf{development}} & \multicolumn{4}{c}{\textbf{test}} \\
\cmidrule(lr){3-6}\cmidrule(lr){7-10}
\textbf{Method} & \textbf{Control} & QWK & Acc19 & far$\geq$3 & far$\geq$5 & QWK & Acc19 & far$\geq$3 & far$\geq$5 \\
\midrule
CORN & CORAL & -0.98 & +20.33 & +0.97 & +1.01 & -1.76 & +20.24 & +1.28 & +1.58 \\
RANK & Regression (MSE) & -0.21 & +2.27 & +0.10 & +0.40 & +0.41 & +2.70 & -1.00 & +0.43 \\
PLR & Regression (MSE) & -0.28 & -4.84 & +0.86 & -0.33 & +0.04 & -4.01 & -0.33 & -0.34 \\
\bottomrule
\end{tabular}
\end{adjustbox}
\end{table}


\subsection{Margin Ranking}

Adding the margin ranking term to the regression loss leaves QWK where it was: the difference to
the regression run is \nlDeltaDevQwkRank\ points on the development set and
\nlDeltaTestQwkRank\ on the test set. On the test set that places RANK, at \nlTestQwkRank, in the
leading group of Section~\ref{sec:res-main}, beside ordinal log loss at \testQwkOLL; on the
development set, at \nlDevQwkRank, it sits a few tenths below it. What the term changes is the
error profile on the test set. Exact accuracy rises by \nlDeltaTestAccRank\ points and the share
of errors of three levels or more falls by \nlAbsDeltaTestFarThreeRank\ points to
\nlTestFarThreeRank\ per cent, the lowest value outside the CORAL family. The hinge only acts on
pairs whose scores are in the wrong order or closer than the margin, which is where a regression
that compresses the scale towards its centre makes its large errors. The development set does not
show the same reduction, however (\nlDeltaDevFarThreeRank\ points), so on one seed the effect is
an indication rather than a measurement.

\subsection{Pairwise Logistic}

The pairwise logistic term also leaves QWK unchanged, at \nlDeltaDevQwkPlr\ points on development
and \nlDeltaTestQwkPlr\ on test, but its effect on the errors is consistent across both splits:
PLR has the fewest errors of five levels or more of every run in this thesis,
\nlDevFarFivePlr\ per cent on development and \nlTestFarFivePlr\ per cent on test, below the
regression method, which held that position before. It pays for this with exact accuracy, which
falls by \nlAbsDeltaTestAccPlr\ points on the test set, and with the lowest macro-F1 of the
table, \nlTestMfPlr. Unlike the hinge, the logistic loss never reaches zero, so it keeps pushing
correctly ordered pairs further apart; the result is the most cautious model of the study, which
rarely makes a large mistake and rarely commits to an exact level. The comparison between RANK
and PLR, which differ only in the pair loss, shows that the choice between a hinge and a logistic
function is not neutral: at the same QWK, one buys exact accuracy and the other buys the tail.

\subsection{CORN}

CORN is the loss the ZAI system paper names as future work \cite{nazzal2025zai}, and its matched
control is the CORAL head. It does not reach CORAL's QWK: it is \nlAbsDeltaDevQwkCorn\ points lower
on development and \nlAbsDeltaTestQwkCorn\ lower on test. It changes the behaviour of the head
completely, though. Its exact accuracy is \nlAbsDeltaTestAccCorn\ points higher than CORAL's on
the test set, at \nlTestAccCorn\ per cent, and its macro-F1 is \nlDeltaTestMfCorn\ points higher,
so it uses the whole scale; its mean distance on the test set, \nlTestDistCorn\ levels, is the
lowest of any run in this thesis; and its share of errors of five levels or more is higher than
CORAL's by \nlDeltaTestFarFiveCorn\ points. On the development set it has the highest QWK of all
the methods whose exact accuracy is above 50 per cent, and on the test set it is second among
them, behind EMD at \testQwkEMD.

This is the profile of the softmax heads of Section~\ref{sec:res-tradeoff}, not of CORAL, and it
explains where that trade-off comes from. With a shared projection, CORAL computes one score per
sentence and compares it with 18 thresholds, which is a regression with learned cut-points, and
it behaves like one. Without the shared projection, each of CORN's 18 questions is a classifier of
its own, and the head behaves like a classifier: it commits to levels, wins exact matches and
occasionally commits far from the gold level. The shared projection, not the binary
decomposition, is what gives CORAL its high QWK and its low exact accuracy.

\subsection{Selection Rule and Summary}

For CORN and RANK the epoch with the lowest development loss is also the epoch with the highest
development QWK, so the selection rule of Section~\ref{sec:res-selection} changes nothing
(\nlQwkselTestGainCorn\ and \nlQwkselTestGainRank\ test QWK points). For PLR the QWK-selected
epoch is one later and gives \nlQwkselTestGainPlr\ points, so the rule does not favour any of the
three.

In summary, the losses brought from readability research in other languages carry over to
Arabic in the sense that matters here: the two ranking losses stay within a few tenths of a QWK
point of their control, the margin ranking loss joins the leading group on the test set, and each
of the three losses moves the error profile in a specific and interpretable direction. None of them beats the leading group on both splits, and with one seed per method the
differences in QWK of a few tenths of a point are not measurements. The robust findings are the
qualitative ones: the pairwise logistic loss is the most effective way found in this thesis to
avoid errors of five levels or more, and removing CORAL's shared projection turns an ordinal head
into a classifier.

\section{Frozen-Backbone Track}
\label{sec:res-frozen}

Table~\ref{tab:res-cpu} reports the second track, in which the encoder is frozen and only a
linear head is trained on cached features, over five seeds.

% generated by gen_numbers.py
\begin{table}[H]
\centering
\footnotesize
\caption{Frozen-backbone track: linear heads on cached AraBERTv2 features, mean $\pm$ standard deviation over five seeds, development set. These numbers are not comparable with the fine-tuning tables.}
\label{tab:res-cpu}
\begin{adjustbox}{max width=\textwidth}
\begin{tabular}{lrrrrrrrr}
\toprule
& \multicolumn{4}{c}{\textbf{pooler features}} & \multicolumn{4}{c}{\textbf{mean-pooled features}} \\
\cmidrule(lr){2-5}\cmidrule(lr){6-9}
\textbf{Method} & QWK & Acc19 & Dist & far$\geq$3 & QWK & Acc19 & Dist & far$\geq$3 \\
\midrule
CE & 71.52\,$\pm$\,0.14 & 40.61 & 1.593 & 22.04 & 74.44\,$\pm$\,0.00 & 46.34 & 1.447 & 20.32 \\
Regression (MSE) & 70.16\,$\pm$\,0.13 & 19.76 & 1.723 & 24.32 & 75.57\,$\pm$\,0.12 & 21.17 & 1.591 & 21.32 \\
OLL ($\alpha=1.5$) & 63.39\,$\pm$\,2.53 & 17.84 & 1.796 & 26.24 & 74.18\,$\pm$\,0.04 & 29.27 & 1.504 & 21.41 \\
EMD & 72.84\,$\pm$\,0.32 & 38.04 & 1.597 & 21.84 & 75.40\,$\pm$\,0.36 & 43.77 & 1.445 & 19.78 \\
SOFT ($\beta=2$) & 72.22\,$\pm$\,0.13 & 40.55 & 1.578 & 21.87 & 75.09\,$\pm$\,0.07 & 46.40 & 1.420 & 19.39 \\
SOFT ($\beta=3$) & 71.93\,$\pm$\,0.13 & 40.68 & 1.582 & 21.89 & 74.89\,$\pm$\,0.08 & 46.43 & 1.429 & 19.76 \\
SOFT ($\beta=4$) & 71.77\,$\pm$\,0.15 & 40.67 & 1.586 & 21.92 & 74.68\,$\pm$\,0.13 & 46.44 & 1.438 & 20.09 \\
SOFT ($\beta=5$) & 71.60\,$\pm$\,0.12 & 40.64 & 1.591 & 22.01 & 74.53\,$\pm$\,0.14 & 46.40 & 1.443 & 20.22 \\
CORAL & 71.03\,$\pm$\,0.17 & 25.14 & 1.687 & 23.34 & 75.53\,$\pm$\,0.16 & 27.47 & 1.559 & 20.24 \\
\midrule
MSE $\rightarrow$ Soft-QWK & 74.55\,$\pm$\,0.06 & 17.90 & 1.838 & 26.85 & 78.00\,$\pm$\,0.10 & 19.68 & 1.702 & 23.17 \\
\midrule
CORAL, zero-initialised & 71.05\,$\pm$\,0.18 & 25.15 & 1.686 & 23.33 & 75.50\,$\pm$\,0.06 & 27.52 & 1.559 & 20.32 \\
\bottomrule
\end{tabular}
\end{adjustbox}
\end{table}


Three results deserve comment.

The first is the variance. Every method except one has a standard deviation over five seeds of
at most \cpuStdMaxQwk\ QWK points, so on frozen features the seed is nearly irrelevant. The
exception is OLL on pooled features, at \cpuStdOll\ points, which makes the point that the
stability of a loss is itself a property worth reporting: the same loss that is among the
leaders when the encoder can adapt is the least stable one when it cannot.

The second is the ranking. Under freezing, the metric-aligned two-phase method is the strongest
on QWK on both representations, \cpuPoolerdevQwkSoftqwk\ and \cpuMeandevQwkSoftqwk, and it
remains the weakest on exact accuracy and mean distance, so the trade-off of
Section~\ref{sec:res-tradeoff} survives the change of regime. The ordering of the remaining
methods does not survive. The rank correlation between the fine-tuned ranking and the frozen
ranking is \cpuSpearmanPooler\ on pooled features and \cpuSpearmanMean\ on mean-pooled
features, that is, no relationship at all in the second case and a mild inversion in the first.
OLL, third when fine-tuned, is last under freezing on both representations, and cross-entropy
climbs. A loss comparison on frozen features is therefore not a cheap substitute for a loss
comparison under fine-tuning, and this thesis reports the frozen track for its variance and for
its CPU-friendly setting rather than as evidence about the fine-tuning ranking.

The third is the cut-point ablation that the fine-tuning track could not afford. With the head
trained for \cpuEpochs\ epochs at $\cpuHeadLr$, the travel budget is \cpuLrSum, and the
frequency initialisation of Equation~\ref{eq:coral-init} then makes no difference at all:
the frequency-initialised and the zero-initialised heads differ by \cpuZeroInitDiffPooler\ QWK
points on pooled features and \cpuZeroInitDiffMean\ on mean-pooled features, which is inside
the seed-to-seed spread. This is the other side of
the argument in Section~\ref{sec:res-cutpoints}. When the budget is large the initialisation is
irrelevant; when it is \lrSumEncoder, as under the BAREC protocol, the cut-points never leave
their initial values. The initialisation is a remedy for a specific arithmetic situation, not a
general improvement, and the frozen track is what shows the difference.

\section{Context: Published Systems and Human Agreement}
\label{sec:res-context}

Two comparisons frame these numbers, and both need a caution.

The published shared-task systems report QWK on the blind test set of the BAREC 2025 shared
task: 85.5 for the CORAL system and 85.6 for the two-phase system
\cite{nazzal2025zai,trigui2025stbw}. Those are different data from the corpus test split used
here, and both systems differ from this study in more than the loss: the CORAL system selects
on QWK, tunes its threshold and uses early stopping, and the two-phase system adds metadata
features. The results here are therefore comparable with the benchmark tables, which they
reproduce, and not with the leaderboard.

Human agreement is the second frame. The corpus authors report an average pairwise QWK between
annotators of \iaaPairwiseQwk\ and an average QWK of \iaaUnifiedQwk\ between individual
annotators and the unified label, with an exact pairwise agreement of only
\iaaPairwiseAcc\ per cent and an average pairwise distance of 0.94 levels
\cite{elmadani2025barec}. The best test QWK in this study, \testQwkCoralHeadlr, is therefore
above the agreement of two randomly chosen annotators and still below the agreement of an
annotator with the adjudicated label. The exact-accuracy figures deserve the same framing: the
best exact accuracy here, \testAccSoftFive\ per cent, is not far below what two annotators
achieve with each other. This is the honest measure of how much room is left on this task, and
it also explains why nineteen-way exact accuracy is a demanding criterion for a scale whose own
annotators agree exactly in six cases out of ten.

\section{Limitations}
\label{sec:res-limits}

The following limitations bound every claim made above.

\begin{enumerate}
\item \textbf{One seed per method in the fine-tuning track.} The protocol fixes the seed at 42
and the accelerator budget did not allow repeats. Differences of a few tenths of a QWK point,
which is the size of the gaps among the four leaders, cannot be separated from run-to-run
variation. This is why no single best method is named. The frozen track shows that seed
variation is small once the encoder is fixed, but it says nothing about variation in
fine-tuning, where the encoder itself moves.

\item \textbf{Mixed precision.} The runs use fp16 autocast where the benchmark runs in single
precision. The losses are computed in single precision, but the encoder forward and backward
passes are not, and the number of optimiser steps skipped because of gradient overflow was not
recorded. A small part of the differences in Table~\ref{tab:res-repro} may come from this.

\item \textbf{Batch composition.} The effective batch of 64 is assembled from four micro-batches
of 16. This is arithmetically equivalent for per-sentence losses except in the last step of each
epoch, as described in Section~\ref{sec:meth-protocol}.

\item \textbf{The two-phase method is not in the comparable group.} It uses the same six epochs
but a different schedule and, in its second phase, a true batch of 64 with gradient
checkpointing. Its numbers answer whether the metric-aligned objective helps, not whether it
would win a comparison in which every method had a two-phase schedule.

\item \textbf{The CORAL cut-point result is a measurement plus one ablation.} That the
cut-points barely move is measured directly and is not in doubt. That the initialisation is
therefore necessary is an argument supported by the movement, by the head learning rate
ablation and by the frozen-track comparison; a zero-initialised CORAL run under the fine-tuning
protocol was not performed and would be the direct test.

\item \textbf{The regression gap.} Section~\ref{sec:res-repro} attributes most of it to the
selected epoch and supports that with the per-epoch values, but the reason our development loss
for that method bottoms out two epochs earlier than the published run does remains unexplained.

\item \textbf{The frozen track is not comparable with the published tables} and uses head
hyper-parameters chosen by us, as stated in Section~\ref{sec:meth-frozen}.

\item \textbf{The added losses.} CORN, RANK and PLR were each trained once, after the main
comparison, and two of their settings are ours: the ranking pairs are formed inside the training
batch, because BAREC has no parallel versions of a text, and the weight of the ranking term is
fixed at one without tuning. Their comparisons with the matched controls share the single-seed
limitation above.

\item \textbf{Scope.} One encoder, one corpus, sentence level, and no ensembling, no external
data and no metadata features. Nothing here speaks to document-level readability or to other
Arabic encoders.
\end{enumerate}

\section{Summary}

The pipeline reproduces the published benchmark: seven of eight losses land within 1.5 QWK
points of their published values and the eighth is explained by one checkpoint decision, which
the per-epoch values in Table~\ref{tab:res-epochs} make visible.

On the question the research brief asks, the answer has three parts. Ordinal-aware objectives
beat the ordinal-blind baseline by up to three QWK points on development and four on test, and
the CORAL head and the metric-aligned two-phase objective are both in the leading group.
Within that group, made up of CORAL, the two-phase Soft-QWK schedule, ordinal log loss and
plain regression, the differences are a few tenths of a point on one seed and no single winner
can be declared. And the ordinal mechanisms buy their QWK specifically by removing far misses,
not by reducing the average error: CORAL has the fewest errors of three levels or more on both
splits, the regression method the fewest of five levels or more, while the lowest mean distance
and the highest exact accuracy belong to a soft-label configuration two QWK points behind them.

Three further findings came out of the ablations. The cut-points of a CORAL head move less than
\cutMoveFineCoral\ under the BAREC learning rate, so they are effectively supplied by their
initialisation rather than learned, and raising only the head learning rate moves them and
improves both QWK and exact accuracy. The checkpoint selection rule is worth up to
\qwkselTestGainMax\ test QWK points and is worth most to the methods whose loss is least
aligned with the metric, so it has to be held fixed in any comparison. And the ranking of
losses does not survive freezing the encoder, which means a cheap frozen-feature study cannot
stand in for a fine-tuning study, even though it is the setting in which seed variance can
actually be measured.

The three losses added from readability research in other languages confirm this picture from
a different direction. The pairwise margin ranking loss joins the leading group on the test set
at \nlTestQwkRank; the pairwise logistic loss has the fewest errors of five levels or more of any
run, \nlTestFarFivePlr\ per cent on test; and CORN, which removes CORAL's shared projection,
gives up \nlAbsDeltaTestQwkCorn\ test QWK points for \nlAbsDeltaTestAccCorn\ points of exact
accuracy and the lowest mean distance of the study, which shows that the shared projection is what
makes CORAL behave like a regression.

\textit{The next chapter draws these findings together and states what remains to be done.}
